% Compilar desde este directorio con:
% pdflatex Clase_01_Introduccion_a_la_Ciencia_de_Datos.tex
\documentclass[aspectratio=169,11pt]{beamer}

\usepackage[utf8]{inputenc}
\usepackage[T1]{fontenc}
\usepackage[spanish,es-nodecimaldot]{babel}
\usepackage{amsmath}
\usepackage{amssymb}
\usepackage{booktabs}
\usepackage{tabularx}
\usepackage{array}
\usepackage{tikz}
\usetikzlibrary{arrows.meta,fit,positioning,shapes.geometric}

\definecolor{UAPNavy}{HTML}{17324D}
\definecolor{UAPTeal}{HTML}{007F82}
\definecolor{UAPGold}{HTML}{E2A33A}
\definecolor{UAPLight}{HTML}{F3F6F8}
\definecolor{UAPGray}{HTML}{536471}

\tikzset{
  flow/.style={-{Stealth[length=2.4mm]},thick,draw=UAPNavy},
  feedback/.style={-{Stealth[length=2.4mm]},thick,dashed,draw=UAPTeal},
  datum/.style={draw=UAPGold,fill=UAPGold!16,rounded corners=2pt,align=center,minimum height=8mm},
  process/.style={draw=UAPTeal,fill=UAPTeal!12,rounded corners=3pt,very thick,align=center,minimum height=10mm},
  decision/.style={draw=UAPNavy,fill=UAPNavy!9,rounded corners=3pt,very thick,align=center,minimum height=10mm},
  human/.style={draw=UAPGold!90!black,fill=UAPGold!18,rounded corners=3pt,very thick,align=center,minimum height=10mm},
  card/.style={draw=UAPTeal,fill=UAPLight,rounded corners=3pt,align=center,minimum height=11mm},
  axiscard/.style={draw=UAPNavy!65,fill=UAPLight,rounded corners=2pt,align=center,minimum height=9mm}
}

\usetheme{Boadilla}
\usecolortheme{default}
\setbeamercolor{normal text}{fg=UAPNavy,bg=white}
\setbeamercolor{structure}{fg=UAPTeal}
\setbeamercolor{frametitle}{fg=white,bg=UAPNavy}
\setbeamercolor{title}{fg=white}
\setbeamercolor{subtitle}{fg=UAPGold}
\setbeamercolor{block title}{fg=white,bg=UAPTeal}
\setbeamercolor{block body}{fg=UAPNavy,bg=UAPLight}
\setbeamercolor{alerted text}{fg=UAPTeal}
\setbeamerfont{frametitle}{series=\bfseries}
\setbeamertemplate{navigation symbols}{}
\setbeamertemplate{itemize item}{\color{UAPGold}$\blacktriangleright$}
\setbeamertemplate{itemize subitem}{\color{UAPTeal}$\bullet$}
\setbeamertemplate{footline}{%
  \leavevmode%
  \hbox{%
    \begin{beamercolorbox}[wd=.78\paperwidth,ht=2.8ex,dp=1.2ex,leftskip=1em]{author in head/foot}%
      \color{UAPGray}\insertshorttitle
    \end{beamercolorbox}%
    \begin{beamercolorbox}[wd=.22\paperwidth,ht=2.8ex,dp=1.2ex,rightskip=1em plus1fil]{date in head/foot}%
      \color{UAPGray}\hfill Semana 1 \enspace | \enspace \insertframenumber/\inserttotalframenumber
    \end{beamercolorbox}%
  }%
}

\newcommand{\concepto}[1]{\textcolor{UAPTeal}{\textbf{#1}}}
\newcommand{\pregunta}[1]{\begin{block}{Pregunta clave}#1\end{block}}
\newcommand{\br}{\\}

\title[Data Science · Clase 1]{Introducción a la Ciencia de Datos}
\subtitle{De una necesidad a evidencia útil para decidir}
\author{Curso de Data Science}
\institute{Semana 1 · Clase 1}
\date{}

\begin{document}

% 1
\begin{frame}[plain]
  \begin{beamercolorbox}[wd=\paperwidth,ht=\paperheight,sep=2em]{frametitle}
    \vfill
    {\usebeamerfont{title}\usebeamercolor[fg]{title}\Huge\inserttitle\par}
    \vspace{0.5em}
    {\usebeamerfont{subtitle}\usebeamercolor[fg]{subtitle}\Large\insertsubtitle\par}
    \vspace{2em}
    {\color{white}\large\insertauthor\par}
    {\color{white}\insertinstitute\par}
    \vfill
  \end{beamercolorbox}
\end{frame}

% 2
\begin{frame}{Propósito de la clase}
  Comprender la Ciencia de Datos como un proceso interdisciplinario que transforma observaciones en evidencia para apoyar decisiones responsables.
  \vspace{0.5em}
  \begin{block}{Al finalizar, podremos}
    \begin{itemize}
      \item distinguir dato, representación, modelo, evidencia y decisión;
      \item diferenciar descripción, predicción y prescripción;
      \item formular unidad, población, alcance y horizonte antes de modelar;
      \item explicar workflow, KDD, CRISP-DM, procedencia y reproducibilidad;
      \item especificar un problema de movilidad con unidad \textbf{zona-franja}.
    \end{itemize}
  \end{block}
\end{frame}

% 3
\begin{frame}{Pregunta de apertura}
  \begin{block}{Situación}
    Una empresa de movilidad quiere ``mejorar el servicio''.
  \end{block}
  \pregunta{¿Qué debemos precisar antes de abrir un dataset?}
  \begin{columns}[T]
    \column{0.48\textwidth}
      \begin{itemize}
        \item ¿Qué situación se desea modificar?
        \item ¿Quién utilizará el resultado?
        \item ¿Sobre qué unidad se decidirá?
      \end{itemize}
    \column{0.48\textwidth}
      \begin{itemize}
        \item ¿Qué acción está disponible?
        \item ¿Qué datos existirán a tiempo?
        \item ¿Cómo se medirá el valor?
      \end{itemize}
  \end{columns}
\end{frame}

% 4
\begin{frame}{¿Qué es la Ciencia de Datos?}
  Campo interdisciplinario que desarrolla métodos para:
  \begin{enumerate}
    \item obtener conocimiento a partir de datos;
    \item comunicar evidencia sobre un fenómeno;
    \item apoyar decisiones bajo objetivos y restricciones.
  \end{enumerate}
  \begin{columns}[T]
    \column{0.48\textwidth}
      \begin{block}{No se reduce a}
        Programar, graficar, entrenar modelos o acumular datos.
      \end{block}
    \column{0.48\textwidth}
      \begin{block}{Producto apropiado}
        Puede ser un inventario, auditoría, análisis, estimación, visualización o recomendación.
      \end{block}
  \end{columns}
  \centering\alert{La pregunta y la decisión determinan el producto.}
\end{frame}

% 5
\begin{frame}{Los datos representan, no reproducen}
  \begin{columns}[c]
    \column{0.58\textwidth}
      \centering
      \begin{tikzpicture}[node distance=7mm,scale=0.86,transform shape]
        \node[decision,minimum width=4.4cm] (phen) {Fenómeno\br demoras en la ciudad};
        \node[process,minimum width=4.4cm,below=of phen] (obs) {Proceso de observación\br sensor · regla · canal};
        \node[datum,minimum width=4.4cm,below=of obs] (rec) {Registro\br hora · GPS · reclamo};
        \draw[flow] (phen) -- (obs);
        \draw[flow] (obs) -- (rec);
        \node[card,right=8mm of obs,minimum width=3.0cm] (out) {Queda fuera\br contexto · silencios\br error · selección};
        \draw[feedback] (obs) -- (out);
      \end{tikzpicture}
    \column{0.38\textwidth}
      Un dato es un registro producido por un proceso situado.
      \vspace{0.7em}
      \begin{block}{Preguntas mínimas}
        ¿Qué representa?\par
        ¿Cómo se obtuvo?\par
        ¿Qué no observa?\par
        ¿Cuándo estuvo disponible?
      \end{block}
  \end{columns}
\end{frame}

% 6
\begin{frame}{De datos a conocimiento}
  \centering
  \begin{tikzpicture}[node distance=5mm,scale=0.88,transform shape]
    \node[datum,minimum width=2.8cm] (d) {\textbf{Datos}\br observaciones\br codificadas};
    \node[process,minimum width=2.8cm,right=of d] (i) {\textbf{Información}\br datos organizados\br en contexto};
    \node[process,minimum width=2.8cm,right=of i] (k) {\textbf{Conocimiento}\br patrones + reglas\br de uso};
    \node[human,minimum width=2.8cm,right=of k] (a) {\textbf{Decisión}\br selección de\br una acción};
    \draw[flow] (d) -- node[above,font=\small]{organizar} (i);
    \draw[flow] (i) -- node[above,font=\small]{interpretar} (k);
    \draw[flow] (k) -- node[above,font=\small]{deliberar} (a);
  \end{tikzpicture}
  \vfill
  \begin{block}{Trazabilidad}
    Cada transición agrega contexto, pero también supuestos, pérdidas y posibles sesgos. Debemos poder reconstruirla.
  \end{block}
\end{frame}

% 7
\begin{frame}{De la necesidad a la decisión}
  \centering
  \begin{tikzpicture}[node distance=4mm,scale=0.79,transform shape]
    \node[decision,minimum width=2.0cm] (need) {Necesidad};
    \node[process,minimum width=2.0cm,right=of need] (q) {Pregunta};
    \node[datum,minimum width=2.0cm,right=of q] (data) {Datos};
    \node[process,minimum width=2.0cm,right=of data] (ana) {Análisis};
    \node[datum,minimum width=2.0cm,right=of ana] (ev) {Evidencia};
    \node[human,minimum width=2.0cm,right=of ev] (dec) {Decisión\br humana};
    \node[decision,minimum width=2.0cm,right=of dec] (imp) {Impacto};
    \foreach \x/\y in {need/q,q/data,data/ana,ana/ev,ev/dec,dec/imp}{\draw[flow] (\x) -- (\y);}
  \end{tikzpicture}
  \vspace{0.8em}
  \begin{columns}[T]
    \column{0.48\textwidth}
      \begin{block}{Correspondencia hacia adelante}
        La tarea debe producir evidencia utilizable por una acción real.
      \end{block}
    \column{0.48\textwidth}
      \begin{block}{Correspondencia hacia atrás}
        La decisión define unidad, oportunidad, evaluación y datos necesarios.
      \end{block}
  \end{columns}
  \centering\alert{Un modelo es un componente de la cadena, no la decisión.}
\end{frame}

% 8
\begin{frame}{El ciclo aprende y vuelve atrás}
  \begin{columns}[c]
    \column{0.58\textwidth}
      \centering
      \begin{tikzpicture}[scale=0.84,transform shape]
        \node[process] (form) at (0,2.0) {Formular};
        \node[datum] (observe) at (3.0,0.7) {Obtener\br y representar};
        \node[process] (analyze) at (1.9,-1.8) {Analizar\br y evaluar};
        \node[human] (use) at (-1.9,-1.8) {Revisar\br y decidir};
        \node[decision] (impact) at (-3.0,0.7) {Observar\br impacto};
        \draw[flow] (form) -- (observe);
        \draw[flow] (observe) -- (analyze);
        \draw[flow] (analyze) -- (use);
        \draw[flow] (use) -- (impact);
        \draw[flow] (impact) -- (form);
        \draw[feedback] (analyze) to[bend right=20] (form);
        \draw[feedback] (use) to[bend left=18] (observe);
      \end{tikzpicture}
    \column{0.38\textwidth}
      \begin{itemize}
        \item Explorar puede revelar que la pregunta no es medible.
        \item Evaluar puede mostrar que no se supera el baseline.
        \item Usar puede cambiar el proceso que genera datos.
      \end{itemize}
      \begin{alertblock}{Autoridad}
        La organización define quién aprueba, ejecuta, revisa y puede detener el sistema.
      \end{alertblock}
  \end{columns}
\end{frame}

% 9
\begin{frame}{Taxonomía de preguntas analíticas}
  \small
  \begin{tabularx}{\textwidth}{>{\bfseries}p{0.18\textwidth} X X}
    \toprule
    Objetivo & Pregunta & Ejemplo en movilidad \\
    \midrule
    Describir & ¿Qué ocurrió? & ¿Dónde y cuándo hubo demoras? \\
    Diagnosticar & ¿Qué se asocia? & ¿Qué acompaña demoras mayores? \\
    Predecir & ¿Qué valor no conocemos? & ¿Qué demanda tendrá cada zona-franja? \\
    Prescribir & ¿Qué acción conviene? & ¿Cómo reasignar vehículos factibles? \\
    Monitorear & ¿Qué está cambiando? & ¿Se degradó la cobertura o el modelo? \\
    \bottomrule
  \end{tabularx}
  \vfill
  \centering\alert{Pueden conectarse, pero requieren evidencia y evaluación diferentes.}
\end{frame}

% 10
\begin{frame}{Describir: resumir lo observado}
  Una descripción estima propiedades de observaciones y grupos:
  \[
    \bar{x}=\frac{1}{n}\sum_{i=1}^{n}x_i
    \qquad
    \widehat{p}=\frac{1}{n}\sum_{i=1}^{n}\mathbb{1}(y_i=1)
  \]
  \begin{columns}[T]
    \column{0.48\textwidth}
      \begin{block}{Productos}
        Frecuencias, medias, cuantiles, distribuciones, segmentos y visualizaciones.
      \end{block}
    \column{0.48\textwidth}
      \begin{block}{Movilidad}
        Viajes, demora mediana y espera extrema por zona y franja observada.
      \end{block}
  \end{columns}
  \centering\alert{Asociación descriptiva no demuestra causa ni anticipa por sí sola.}
\end{frame}

% 11
\begin{frame}{Predecir: estimar lo no observado}
  Una predicción estima un resultado desconocido, futuro o costoso de medir:
  \[
    \widehat{P}(Y\mid X)
    \qquad\text{o}\qquad
    \widehat{\mathbb{E}}[Y\mid X]
  \]
  \begin{itemize}
    \item Las entradas deben existir en el momento real de uso.
    \item La evaluación utiliza casos no usados para ajustar.
    \item La partición debe respetar tiempo, grupos y duplicados.
    \item El desempeño se compara con un baseline pertinente.
  \end{itemize}
  \pregunta{¿Qué demanda tendrá cada zona durante la próxima franja?}
\end{frame}

% 12
\begin{frame}{Describir no es predecir}
  \centering
  \begin{tikzpicture}[node distance=11mm,scale=0.90,transform shape]
    \node[card,minimum width=5.0cm,minimum height=2.0cm] (desc) {\textbf{Descripción}\br[1mm]resume casos observados\br $\bar{x}$ · cuantiles · distribución};
    \node[card,minimum width=5.0cm,minimum height=2.0cm,right=of desc] (pred) {\textbf{Predicción}\br[1mm]estima casos no observados\br $\widehat{P}(Y\mid X)$ · error futuro};
    \draw[flow] (desc) -- node[above,font=\small]{hipótesis, no garantía} (pred);
    \node[datum,below=8mm of desc,minimum width=5.0cm] {¿Dónde ocurrieron demoras?};
    \node[datum,below=8mm of pred,minimum width=5.0cm] {¿Qué zona-franja tendrá demanda alta?};
  \end{tikzpicture}
  \vfill
  \begin{block}{Prueba decisiva}
    Una buena explicación del pasado no implica bajo error sobre periodos posteriores.
  \end{block}
\end{frame}

% 13
\begin{frame}{Predecir no es decidir}
  \begin{columns}[c]
    \column{0.61\textwidth}
      \centering
      \begin{tikzpicture}[node distance=6mm,scale=0.86,transform shape]
        \node[datum,minimum width=3.0cm] (pred) {Predicción\br demanda +\br incertidumbre};
        \node[process,minimum width=3.0cm,right=of pred] (compare) {Comparar acciones\br costos · restricciones\br consecuencias};
        \node[human,minimum width=3.0cm,right=of compare] (person) {Responsable\br aprueba, modifica\br o rechaza};
        \draw[flow] (pred) -- (compare);
        \draw[flow] (compare) -- (person);
      \end{tikzpicture}
    \column{0.35\textwidth}
      \[
        a^*=\arg\max_{a\in\mathcal A}
        \mathbb E[U(a,S)\mid X]
      \]
      \begin{itemize}
        \item El modelo informa.
        \item La política compara.
        \item La persona autorizada decide.
      \end{itemize}
  \end{columns}
  \begin{alertblock}{Autoridad humana efectiva}
    Requiere tiempo, información y poder real para contradecir o detener la recomendación.
  \end{alertblock}
\end{frame}

% 14
\begin{frame}{Una disciplina interdisciplinaria}
  \centering
  \begin{tikzpicture}[scale=0.81,transform shape]
    \node[process,circle,minimum size=2.6cm] (ds) at (0,0) {\textbf{Evidencia}\br útil y\br confiable};
    \node[card,minimum width=3.2cm] (stat) at (-4.0,1.8) {Estadística\br variabilidad e inferencia};
    \node[card,minimum width=3.2cm] (math) at (0,2.8) {Matemática\br objetivos y restricciones};
    \node[card,minimum width=3.2cm] (comp) at (4.0,1.8) {Computación\br representación y escala};
    \node[card,minimum width=3.2cm] (eng) at (4.0,-1.8) {Ingeniería\br pruebas y operación};
    \node[card,minimum width=3.2cm] (domain) at (0,-2.8) {Dominio\br significado y plausibilidad};
    \node[card,minimum width=3.2cm] (comm) at (-4.0,-1.8) {Comunicación\br comprensión y uso};
    \foreach \x in {stat,math,comp,eng,domain,comm}{\draw[flow] (\x) -- (ds);}
  \end{tikzpicture}
  \vfill
  \centering\alert{Una perspectiva aislada puede optimizar la parte equivocada.}
\end{frame}

% 15
\begin{frame}{Formular antes de modelar}
  \centering
  \begin{tikzpicture}[node distance=6mm,scale=0.90,transform shape]
    \node[decision,minimum width=3.6cm] (need) {\textbf{Necesidad real}\br reducir esperas};
    \node[process,minimum width=4.0cm,right=of need] (question) {\textbf{Pregunta analítica}\br demanda por zona-franja};
    \node[datum,minimum width=3.8cm,right=of question] (task) {\textbf{Tarea computacional}\br agregar o predecir};
    \draw[flow] (need) -- (question);
    \draw[flow] (question) -- (task);
    \draw[feedback] (task.south) -- ++(0,-7mm) -| node[pos=.25,below,font=\small]{¿sirve para decidir?} (need.south);
  \end{tikzpicture}
  \vspace{0.8em}
  \begin{block}{Prueba sin algoritmos}
    Si no podemos explicar quién decide, sobre qué unidad, con qué información, para cuándo y con qué propósito, la formulación aún no está completa.
  \end{block}
\end{frame}

% 16
\begin{frame}{Especificación compacta del problema}
  \[
    \mathcal P=(U,\mathcal R,O,X,Y,A,H,C)
  \]
  \small
  \begin{columns}[T]
    \column{0.48\textwidth}
      \begin{itemize}
        \item $U$: unidad de análisis.
        \item $\mathcal R$: población objetivo.
        \item $O$: objetivo del proyecto.
        \item $X$: entradas disponibles.
        \item $Y$: resultado de interés.
      \end{itemize}
    \column{0.48\textwidth}
      \begin{itemize}
        \item $A$: acciones posibles.
        \item $H$: horizonte temporal.
        \item $C$: restricciones.
        \item Usuario, valor y responsables completan la ficha.
      \end{itemize}
  \end{columns}
  \begin{block}{Cuatro pruebas}
    Observabilidad · accionabilidad · disponibilidad temporal · contrafactual: ``¿qué haríamos si la respuesta fuera distinta?''
  \end{block}
\end{frame}

% 17
\begin{frame}{Unidad, población y alcance}
  \begin{columns}[T]
    \column{0.32\textwidth}
      \begin{block}{Unidad $U$}
        Entidad elemental descrita, predicha o decidida. Define el significado de una fila.
      \end{block}
    \column{0.32\textwidth}
      \begin{block}{Población $\mathcal R$}
        Unidades sobre las que se desea generalizar o actuar; no equivale a la muestra observada.
      \end{block}
    \column{0.32\textwidth}
      \begin{block}{Alcance}
        Periodos, lugares, casos, usuarios, acciones y exclusiones autorizadas.
      \end{block}
  \end{columns}
  \vspace{0.5em}
  \begin{alertblock}{Coherencia de granularidad}
    Si la acción reasigna recursos entre zonas por franja, la evidencia final debe conservar la unidad \textbf{zona-franja}; predicciones por viaje requieren una agregación explícita.
  \end{alertblock}
  \centering Población objetivo $\neq$ muestra disponible $\neq$ población futura de uso.
\end{frame}

% 18
\begin{frame}{Horizonte y disponibilidad: reconstruir el reloj}
  \centering
  \begin{tikzpicture}[x=1cm,y=1cm,scale=0.90,transform shape]
    \draw[thick,UAPNavy,-{Stealth[length=2.4mm]}] (-6,0) -- (6,0);
    \foreach \x/\lab in {-5/{histórico},-1/{corte $t_0$},2/{franja objetivo},5/{resultado}}{
      \draw[thick,UAPNavy] (\x,-0.12) -- (\x,0.12);
      \node[below] at (\x,-0.18) {\lab};
    }
    \node[datum,minimum width=4.0cm] at (-3,1.2) {Entradas $X$ conocidas\br antes de $t_0$};
    \node[human,minimum width=3.0cm] at (-1,-1.25) {Decisión\br en $t_0$};
    \node[decision,minimum width=3.5cm] at (3.5,1.2) {Objetivo $Y$\br se observa después};
    \draw[flow] (-3,0.75) -- (-1,-0.75);
    \draw[feedback] (3.5,0.75) -- node[above,font=\small]{prohibido como entrada} (-0.7,0.25);
  \end{tikzpicture}
  \vfill
  \begin{block}{Regla temporal}
    Evaluar con la información disponible retrospectivamente puede sobrestimar lo que el sistema sabía al decidir.
  \end{block}
\end{frame}

% 19
\begin{frame}{Workflow de un proyecto de datos}
  \centering
  \begin{tikzpicture}[node distance=4mm,scale=0.78,transform shape]
    \node[decision,minimum width=2.6cm] (f) {1. Formular\br decisión y éxito};
    \node[datum,minimum width=2.6cm,right=of f] (o) {2. Obtener\br e inventariar};
    \node[process,minimum width=2.6cm,right=of o] (p) {3. Preparar\br representación};
    \node[process,minimum width=2.6cm,right=of p] (m) {4. Analizar\br o modelar};
    \node[datum,minimum width=2.6cm,right=of m] (e) {5. Evaluar\br y contrastar};
    \foreach \x/\y in {f/o,o/p,p/m,m/e}{\draw[flow] (\x) -- (\y);}
    \node[human,minimum width=4.0cm,below=12mm of p] (u) {6. Comunicar o integrar\br con autoridad y monitoreo};
    \draw[flow] (e.south) |- (u.east);
    \draw[feedback] (u.west) -| (f.south);
  \end{tikzpicture}
  \vfill
  \begin{block}{Principio}
    Baseline, criterios de éxito y protocolo de evaluación se fijan antes de comparar resultados finales.
  \end{block}
\end{frame}

% 20
\begin{frame}{Puertas y artefactos: aprender antes de escalar}
  \small
  \begin{tabularx}{\textwidth}{>{\bfseries}p{0.21\textwidth} X X}
    \toprule
    Puerta & Evidencia mínima & Artefacto vivo \\
    \midrule
    Formular $\rightarrow$ datos & Unidad, alcance, autoridad y valor acordados & Ficha del problema \\
    Datos $\rightarrow$ análisis & Cobertura, licencia, calidad y tiempo conocidos & Inventario + diccionario \\
    Análisis $\rightarrow$ uso & Baseline superado y riesgos aceptables & Registro de experimento \\
    Uso $\rightarrow$ escala & Utilidad, carga y fallos observados & Plan de monitoreo y retirada \\
    \bottomrule
  \end{tabularx}
  \vfill
  \begin{block}{Producto mínimo evaluable}
    El artefacto más pequeño que permite probar la incertidumbre capaz de invalidar el proyecto.
  \end{block}
  \centering\alert{Continuar, modificar, escalar o detener son resultados legítimos.}
\end{frame}

% 21
\begin{frame}{KDD: del dato al conocimiento validado}
  \centering
  \begin{tikzpicture}[node distance=4mm,scale=0.83,transform shape]
    \node[datum,minimum width=2.25cm] (s) {Selección};
    \node[process,minimum width=2.25cm,right=of s] (p) {Preprocesamiento};
    \node[process,minimum width=2.25cm,right=of p] (t) {Transformación};
    \node[process,minimum width=2.25cm,right=of t] (m) {Minería\br de datos};
    \node[decision,minimum width=2.65cm,right=of m] (i) {Interpretación\br y evaluación};
    \foreach \x/\y in {s/p,p/t,t/m,m/i}{\draw[flow] (\x) -- (\y);}
    \draw[feedback] (i.south) -- ++(0,-7mm) -| node[pos=.25,below,font=\small]{iterar} (p.south);
  \end{tikzpicture}
  \vspace{0.8em}
  \begin{columns}[T]
    \column{0.48\textwidth}
      \begin{block}{Distinción}
        La minería es una etapa dentro de KDD, no un sinónimo del proceso completo.
      \end{block}
    \column{0.48\textwidth}
      \begin{block}{Salto que debe evitarse}
        Resultado $\rightarrow$ patrón $\rightarrow$ hallazgo validado $\rightarrow$ regla aprobada.
      \end{block}
  \end{columns}
\end{frame}

% 22
\begin{frame}{CRISP-DM: organizar el proyecto completo}
  \begin{columns}[c]
    \column{0.60\textwidth}
      \centering
      \begin{tikzpicture}[scale=0.78,transform shape]
        \node[decision] (business) at (0,2.3) {Comprensión\br del problema};
        \node[datum] (data) at (3.2,1.1) {Comprensión\br de datos};
        \node[process] (prep) at (3.2,-1.5) {Preparación};
        \node[process] (model) at (0,-2.6) {Modelado};
        \node[datum] (eval) at (-3.2,-1.5) {Evaluación};
        \node[human] (deploy) at (-3.2,1.1) {Despliegue\br o entrega};
        \draw[flow] (business)--(data);
        \draw[flow] (data)--(prep);
        \draw[flow] (prep)--(model);
        \draw[flow] (model)--(eval);
        \draw[flow] (eval)--(deploy);
        \draw[flow] (deploy)--(business);
        \draw[feedback] (eval) to[bend left=15] (data);
      \end{tikzpicture}
    \column{0.36\textwidth}
      \begin{itemize}
        \item Comienza por el problema, no por el modelo.
        \item Evalúa técnica y correspondencia con el objetivo.
        \item Despliegue puede ser informe, tablero, API o regla operativa.
      \end{itemize}
      \begin{alertblock}{No termina al publicar}
        Operación, monitoreo, incidentes y retirada retroalimentan el ciclo.
      \end{alertblock}
  \end{columns}
\end{frame}

% 23
\begin{frame}{Estructura de los datos}
  \small
  \begin{tabularx}{\textwidth}{>{\bfseries}p{0.20\textwidth} X X}
    \toprule
    Tipo & Organización & Ejemplo y decisión de representación \\
    \midrule
    Estructurados & Esquema explícito, filas, columnas y claves & Viajes SQL; validar unidad, tipos y cardinalidad \\
    Semiestructurados & Claves o marcas con esquema flexible & Eventos JSON; distinguir ausente, nulo y lista vacía \\
    No estructurados & Sin variables tabulares inmediatas & Texto, audio o imagen; construir una representación \\
    \bottomrule
  \end{tabularx}
  \vfill
  \begin{block}{Advertencia semántica}
    Un esquema válido no garantiza significado correcto; y ``no estructurado'' no significa carente de estructura interna.
  \end{block}
\end{frame}

% 24
\begin{frame}{De fuentes heterogéneas a una tabla analítica}
  \centering
  \begin{tikzpicture}[node distance=5mm and 8mm,scale=0.82,transform shape]
    \node[datum,minimum width=2.8cm] (trips) {Viajes\br unidad: viaje};
    \node[datum,minimum width=2.8cm,below=of trips] (weather) {Clima\br estación-hora};
    \node[datum,minimum width=2.8cm,below=of weather] (zones) {Zonas\br polígono-versión};
    \node[process,minimum width=3.6cm,right=of weather] (rules) {Reglas explícitas\br agregar · unir\br filtrar · cortar};
    \node[decision,minimum width=4.1cm,right=of rules] (table) {Tabla analítica\br una fila = zona-franja\br claves + procedencia};
    \draw[flow] (trips) -- (rules);
    \draw[flow] (weather) -- (rules);
    \draw[flow] (zones) -- (rules);
    \draw[flow] (rules) -- (table);
  \end{tikzpicture}
  \vspace{0.6em}
  \begin{itemize}
    \item La tabla analítica deriva de fuentes: no las reemplaza.
    \item Cada unión declara claves, cardinalidad, cobertura y fecha de corte.
    \item Cambiar la cantidad de filas puede cambiar la población o la unidad.
  \end{itemize}
\end{frame}

% 25
\begin{frame}{Fuentes: inventariar antes de integrar}
  \small
  \begin{tabularx}{\textwidth}{>{\bfseries}p{0.16\textwidth} p{0.19\textwidth} p{0.17\textwidth} X}
    \toprule
    Fuente & Unidad original & Acceso & Riesgo principal \\
    \midrule
    Viajes & viaje & base o archivo & operadores no cubiertos \\
    GPS & dispositivo-evento & flujo o API & latencia, deriva y pérdida \\
    Clima & estación-hora & API & desfase horario y distancia \\
    Zonas & polígono-versión & archivo geográfico & límites modificados \\
    Calendario & día & archivo o tabla & definición local \\
    \bottomrule
  \end{tabularx}
  \vfill
  \begin{block}{Ficha mínima por fuente}
    Propietario · propósito original · unidad · cobertura · captura · formato · frecuencia · licencia · sensibilidad · calidad · cambios
  \end{block}
\end{frame}

% 26
\begin{frame}{Procedencia: ejes distintos, preguntas distintas}
  \centering
  \begin{tikzpicture}[node distance=5mm,scale=0.82,transform shape]
    \node[axiscard,minimum width=11.6cm] (gen) {\textbf{Mecanismo de generación}\quad observacional \;|\; experimental \;|\; simulado};
    \node[axiscard,minimum width=11.6cm,below=of gen] (purpose) {\textbf{Propósito de recolección}\quad primario para el proyecto \;|\; secundario, otro propósito};
    \node[axiscard,minimum width=11.6cm,below=of purpose] (access) {\textbf{Acceso y gobernanza}\quad abierto \;|\; privado \;|\; restringido};
    \node[axiscard,minimum width=11.6cm,below=of access] (lineage) {\textbf{Linaje}\quad origen $\rightarrow$ extracción $\rightarrow$ transformación $\rightarrow$ producto};
  \end{tikzpicture}
  \vspace{0.6em}
  \begin{block}{No son categorías excluyentes entre ejes}
    Un dato puede ser observacional, secundario y privado a la vez. Cada eje implica límites de inferencia, uso y auditoría diferentes.
  \end{block}
\end{frame}

% 27
\begin{frame}{Big Data: seis V, seis diagnósticos}
  \centering
  \begin{tikzpicture}[scale=0.82,transform shape]
    \node[process,circle,minimum size=2.3cm] (big) at (0,0) {\textbf{Desafío}\br de datos};
    \node[card,minimum width=3.3cm] (v1) at (-4.3,2.0) {\textbf{Volumen}\br ¿cuánto?};
    \node[card,minimum width=3.3cm] (v2) at (0,2.8) {\textbf{Velocidad}\br ¿con qué latencia?};
    \node[card,minimum width=3.3cm] (v3) at (4.3,2.0) {\textbf{Variedad}\br ¿qué formatos y fuentes?};
    \node[card,minimum width=3.3cm] (v4) at (4.3,-2.0) {\textbf{Veracidad}\br ¿qué confiabilidad?};
    \node[card,minimum width=3.3cm] (v5) at (0,-2.8) {\textbf{Valor}\br ¿mejora la decisión?};
    \node[card,minimum width=3.3cm] (v6) at (-4.3,-2.0) {\textbf{Variabilidad}\br ¿qué cambia?};
    \foreach \x in {v1,v2,v3,v4,v5,v6}{\draw[flow] (\x) -- (big);}
  \end{tikzpicture}
  \vfill
  \centering\alert{Escala no corrige semántica: más datos del concepto equivocado siguen siendo equivocados.}
\end{frame}

% 28
\begin{frame}{Herramientas y reproducibilidad}
  \begin{columns}[T]
    \column{0.31\textwidth}
      \begin{block}{Entornos}
        Python · R · SQL · Jupyter · Quarto
      \end{block}
    \column{0.31\textwidth}
      \begin{block}{Artefactos}
        Código · datos · configuración · resultados · documentación
      \end{block}
    \column{0.31\textwidth}
      \begin{block}{Controles}
        Versiones · checksums · pruebas · semillas · registro de decisiones
      \end{block}
  \end{columns}
  \vspace{0.6em}
  \begin{block}{Cadena reproducible}
    \centering Fuente identificada $\rightarrow$ transformación versionada $\rightarrow$ experimento registrado $\rightarrow$ resultado reconstruible
  \end{block}
  \begin{itemize}
    \item Una semilla ayuda, pero no sustituye datos, entorno y partición versionados.
    \item Un notebook debe ejecutarse de principio a fin en un entorno limpio.
    \item La herramienta se elige por el ciclo de vida y el equipo, no por preferencia personal.
  \end{itemize}
\end{frame}

% 29
\begin{frame}{Caso movilidad: formulación coherente}
  \begin{block}{Decisión operativa}
    Antes de cada franja, el responsable decide si reasigna vehículos entre zonas.
  \end{block}
  \small
  \begin{columns}[T]
    \column{0.48\textwidth}
      \begin{itemize}
        \item \concepto{Unidad}: zona-franja.
        \item \concepto{Población}: zonas operadas en condiciones normales.
        \item \concepto{Horizonte}: próxima franja de 30 minutos.
        \item \concepto{Resultado}: demanda y espera por zona-franja.
      \end{itemize}
    \column{0.48\textwidth}
      \begin{itemize}
        \item \concepto{Acciones}: mantener o proponer traslados factibles.
        \item \concepto{Restricciones}: capacidad, distancia, seguridad y cobertura.
        \item \concepto{Valor}: reducir espera extrema sin aumentar cancelaciones.
        \item \concepto{Autoridad}: responsable de operaciones.
      \end{itemize}
  \end{columns}
  \centering\alert{La unidad coincide con el nivel de predicción y de reasignación.}
\end{frame}

% 30
\begin{frame}{Movilidad: evidencia, recomendación y autoridad}
  \centering
  \begin{tikzpicture}[node distance=5mm,scale=0.80,transform shape]
    \node[datum,minimum width=2.6cm] (sources) {Viajes · GPS\br clima · zonas};
    \node[process,minimum width=2.8cm,right=of sources] (unit) {Construir\br zona-franja};
    \node[datum,minimum width=2.8cm,right=of unit] (evidence) {Demanda esperada\br + incertidumbre};
    \node[process,minimum width=2.8cm,right=of evidence] (options) {Opciones factibles\br + consecuencias};
    \node[human,minimum width=2.8cm,right=of options] (human) {Responsable\br decide};
    \foreach \x/\y in {sources/unit,unit/evidence,evidence/options,options/human}{\draw[flow] (\x) -- (\y);}
    \node[decision,below=10mm of options,minimum width=3.4cm] (fleet) {Acción en la flota\br o mantener estado};
    \draw[flow] (human) |- (fleet);
    \draw[feedback] (fleet.west) -| node[pos=.2,below,font=\small]{impacto observado} (unit.south);
  \end{tikzpicture}
  \vfill
  \begin{alertblock}{Límite de automatización}
    El sistema recomienda y puede abstenerse; no ejecuta traslados ni reemplaza protocolos de seguridad.
  \end{alertblock}
\end{frame}

% 31
\begin{frame}{Del caso amplio a tareas concretas}
  \small
  \begin{tabularx}{\textwidth}{>{\bfseries}p{0.17\textwidth} X X}
    \toprule
    Objetivo & Tarea sobre zona-franja & Producto y evaluación \\
    \midrule
    Describir & Agregar viajes, espera y oferta observada & Perfil por zona-franja \\
    Predecir & Estimar demanda o espera de la próxima franja & Pronóstico vs. baseline temporal \\
    Prescribir & Comparar reasignaciones factibles & Opciones, utilidad y restricciones \\
    Monitorear & Detectar cambio en datos, error e impacto & Alertas, revisión o retirada \\
    \bottomrule
  \end{tabularx}
  \vfill
  \begin{block}{Granularidad}
    Los eventos de viaje se agregan a zona-franja; el resultado no se presenta como una decisión por viaje individual.
  \end{block}
\end{frame}

% 32
\begin{frame}{Fuga temporal: saber hoy lo que mañana ocurrió}
  \centering
  \begin{tikzpicture}[x=1cm,y=1cm,scale=0.86,transform shape]
    \draw[thick,UAPNavy,-{Stealth[length=2.4mm]}] (-6,0)--(6,0);
    \node[human] (cut) at (-1,-1.25) {Decidir 09:00\br para 09:00--09:30};
    \node[datum] (ok1) at (-4,1.2) {GPS recibido\br hasta 08:55};
    \node[datum] (ok2) at (-1.8,2.2) {Calendario y\br flota disponible};
    \node[decision] (bad1) at (2.1,2.0) {Viajes completados\br a las 09:20};
    \node[decision] (bad2) at (4.7,1.0) {Demanda total\br de la franja};
    \draw[flow] (ok1)--(cut);
    \draw[flow] (ok2)--(cut);
    \draw[feedback] (bad1)--node[above,sloped,font=\small]{fuga} (cut);
    \draw[feedback] (bad2)--node[below,sloped,font=\small]{fuga} (cut);
    \draw[thick,UAPGold] (-1,-0.2)--(-1,0.2);
  \end{tikzpicture}
  \vfill
  \begin{block}{Prueba}
    Reconstruir para cada fila qué eventos habían sido recibidos, no solo cuáles ocurrieron antes. Tiempo de evento, recepción y decisión son relojes distintos.
  \end{block}
\end{frame}

% 33
\begin{frame}{Actividad guiada: ficha del problema}
  \small
  En equipos, completar para movilidad:
  \begin{enumerate}
    \item necesidad, usuario, responsable y personas afectadas;
    \item unidad zona-franja, población, alcance y exclusiones;
    \item decisión, acciones permitidas y autoridad final;
    \item fecha de corte, horizonte, entradas disponibles y resultado;
    \item pregunta descriptiva, predictiva y prescriptiva;
    \item fuentes, procedencia y regla de integración;
    \item baseline, valor, restricciones y condición de detención;
    \item un riesgo de fuga temporal y su control.
  \end{enumerate}
  \begin{block}{Tiempo sugerido}
    30 minutos de trabajo + 10 minutos de contraste.
  \end{block}
\end{frame}

% 34
\begin{frame}{Criterios de revisión}
  La ficha es adecuada si:
  \begin{itemize}
    \item la unidad \textbf{zona-franja} coincide con pregunta, tabla y acción;
    \item población, alcance, horizonte y exclusiones son verificables;
    \item cada entrada existe y puede recibirse antes de decidir;
    \item descripción, predicción y prescripción no se confunden;
    \item las fuentes conservan procedencia y reglas de integración;
    \item el baseline representa una alternativa realista;
    \item valor, riesgos y criterios de detención están explícitos;
    \item una persona con autoridad aprueba, modifica o rechaza la acción.
  \end{itemize}
  \begin{block}{Entregable de semana 1}
    Ficha de una página + diagrama necesidad $\rightarrow$ evidencia $\rightarrow$ decisión $\rightarrow$ impacto.
  \end{block}
\end{frame}

% 35
\begin{frame}{Puesta en común}
  Cada equipo presenta en un minuto:
  \begin{enumerate}
    \item la decisión y quién conserva autoridad;
    \item la unidad, el corte temporal y el horizonte;
    \item el baseline y la métrica de valor;
    \item el supuesto que podría invalidar el proyecto.
  \end{enumerate}
  \vspace{0.5em}
  \begin{block}{Preguntas para contrastar}
    \begin{itemize}
      \item ¿La evidencia llega antes y a la misma granularidad que la decisión?
      \item ¿Qué representa la muestra y qué zonas quedan fuera?
      \item ¿Qué resultado justificaría modificar o detener el proyecto?
    \end{itemize}
  \end{block}
\end{frame}

% 36
\begin{frame}{Síntesis: siete ideas para conservar}
  \begin{enumerate}
    \item Los datos son representaciones parciales producidas por procesos.
    \item La Ciencia de Datos construye evidencia, no solo modelos.
    \item Describir, predecir y decidir son problemas distintos.
    \item Unidad, población, alcance y horizonte preceden al algoritmo.
    \item Workflow, KDD y CRISP-DM son iterativos y trazables.
    \item Procedencia y reproducibilidad sostienen lo que podemos afirmar.
    \item La autoridad y responsabilidad pertenecen a personas y organizaciones.
  \end{enumerate}
  \vfill
  \begin{block}{Idea de cierre}
    Primero se diseña la decisión y la evidencia necesaria; después se eligen datos, métodos y herramientas.
  \end{block}
\end{frame}

% 37
\begin{frame}{Síntesis visual: la cadena completa}
  \centering
  \begin{tikzpicture}[node distance=4mm,scale=0.76,transform shape]
    \node[decision,minimum width=2.4cm] (need) {Necesidad\br y alcance};
    \node[process,minimum width=2.4cm,right=of need] (form) {Unidad\br población\br horizonte};
    \node[datum,minimum width=2.4cm,right=of form] (sources) {Fuentes\br y procedencia};
    \node[process,minimum width=2.4cm,right=of sources] (repr) {Representación\br analítica};
    \node[datum,minimum width=2.4cm,right=of repr] (ev) {Evidencia\br evaluada};
    \node[human,minimum width=2.4cm,right=of ev] (dec) {Decisión\br autorizada};
    \foreach \x/\y in {need/form,form/sources,sources/repr,repr/ev,ev/dec}{\draw[flow] (\x)--(\y);}
    \node[decision,minimum width=4.0cm,below=12mm of repr] (impact) {Impacto · monitoreo · aprendizaje};
    \draw[flow] (dec.south) |- (impact.east);
    \draw[feedback] (impact.west) -| (need.south);
    \node[card,minimum width=3.1cm,below=9mm of need] {Pregunta y valor};
    \node[card,minimum width=3.1cm,below=9mm of sources] {Calidad y licencia};
    \node[card,minimum width=3.1cm,below=9mm of ev] {Baseline y riesgos};
  \end{tikzpicture}
  \vfill
  \centering\alert{Cada flecha es una decisión documentable; cada retorno es una oportunidad de aprender.}
\end{frame}

% 38
\begin{frame}{Lecturas y recursos}
  \begin{block}{Lectura principal}
    \begin{itemize}
      \item \textbf{Capítulo 1}: secciones 1.1.1, 1.1.4, 1.1.5 y 1.1.6.
      \item \textbf{Capítulo 2}: secciones 2.1 a 2.4, con énfasis en 2.1.3, 2.2.1, 2.2.2, 2.3 y 2.4.
    \end{itemize}
  \end{block}
  \begin{block}{Para continuar}
    \begin{itemize}
      \item Revisar los glosarios de ambos capítulos.
      \item Completar la ficha inicial del caso de movilidad.
      \item Dibujar el linaje de una fuente hasta la unidad zona-franja.
    \end{itemize}
  \end{block}
\end{frame}

\end{document}
